\section{Legal, Ethical, Social, and Professional Issues}

When submitting the proposal for thus project, a requirement was to consider the associated \acp{lespi}.
The original impacts in these areas where defined in the \ac{ppd} and can be seen in Appendix \ref{app:original-legal-ethical-social-and-professional-issues}.

The following sections cover the \acp{lespi} relating to the now finalised project.

\subsection{Legal Issues}

The core concept of this project is to use \ac{gr} to passively identify individuals.
This relies on the collection and processing of behavioural biometric data, which falls under the "special category" of data under the GDPR and the Data Protection Act 2018 \parencite{gdprinfoGeneralDataProtection2024,legislation.gov.ukDataProtectionAct}. This was outlined in the \ac{ppd} in Section 8.1.

During the devleopment and testing process, appropriate measures were taken to ensure compliance. This included:

\begin{itemize}
    \item Obtaining informed consent from all participants via a participant information sheet and informed consent form.
    \item Anonymising collected data wherever possible (e.g., storing \acp{gei} under generic identifiers and not retaining raw footage).
    \item Limiting data use to the purposes outlined in the participant information sheet.
    \item Disposing of ALL primary data securely after testing concluded.
\end{itemize}

This ensured that research remained ethical and volunteers were not subjected to any unnecessary risks. 
It also ensured all legal obligations were met by the researcher and supervisor during development and testing.

\subsection{Ethical Issues}

Although arguably not as susceptible to biases as other biometric recognition technologies, \ac{gr} is still susceptible to biases relating primarily to differences in physical build, correlating mostly with sex and age.
Primary and secondary datasets used during testing processes included a mix of males and females, with a not-explititly-targeted, but varying nonetheless, age range.
This reduced the risk of bias which, if present, would have invalidated the results and claimed accuracy figures.

Risk of data misuse was another ethical concern, echoed by feedback received during surveying and post-testing questionnaires.
This has been acknowledged during the testing process through participant information sheets and informed consent forms, with all primary data being both anonymised during collection and securely disposed of afterwards.

The potential for nefarious use of \ac{gr} was not explored, with surveillance falling outside of the scope of this project. 

\subsection{Social Issues}

Social considerations relate primarily to the usability of the solution, as well as the potential impact on society.
While \ac{gr}-based authentication enables user-independent authentication, improving convenience, it does not account for individuals with disabilities, as acknowledged in the \ac{ppd} (Section 8.2.2).

Though not addressed in the proposed or finalised solution, multi-modal authentication was identified as a potential solution to this issue, providing alternate methods of authentication for users with disabilities or when \ac{gr} is otherwise unavailable.

As mentioned in the ethical issue assessment, the use of \ac{gr} for surveillance could have a negative societal impact, invading privacy and potentially leading to a 'panopticon' effect --- a concept coined by Michel Foucault in which society is monitored by the state to maintain control \parencite{foucaultm.DisciplinePunishBirth1977}

\subsection{Professional Issues}

Throughout planning, development, and testing, this project adhered to the BCS Code of Conduct. 
Efforts were made to uphold public trust (e.g., informing participants of what data was being collected and how it would be used) and ensure the system was built with security and fairness as priorities, with research conducted with integrity and proper attribution \parencite{bcsBCSCodeConduct2024}.

This project in no way aimed to encourage or enable unethical practices, with the resulting system designed to highlight the potential of \ac{gr} for improving security and convenience, using technology that can unfortunately enable surveillance --- invading privacy.

